\section{Lower Bounds for General Circuits}
\label{sec:demorgan-LB}

We state the following lemma general enough so that we can apply it
for the degree as well as the size lower bound. As explained
previously, for the size lower bounds we rely on functions that almost
have circuits of size $s$. Recall that we consider the class of
functions $\calF_n(s,t)$ that consists of all Boolean functions
$f: \set{0,1}^n \rightarrow \set{0,1}$ for which there is a Boolean
circuit $C_f: \set{0,1}^n \rightarrow \set{0, 1, \bot}$ of size at most
$s$ such that
\begin{enumerate}
\item if $C_f(\alpha) \neq \bot$, then $C_f(\alpha) = f(\alpha)$, and
\item $C_f(\alpha) = \bot$ on at most $t$ inputs.
\end{enumerate}
The following lemma establishes the existence of $m$-independent
$k$-deter\-mined variable substitutions that result in XOR-CSP instances
over explicit graphs.

\begin{lemma}
  \label{lem:restriction}
  For all $k, m, n, t \in \N$ satisfying $m \le 2^n$, and any explicit
  bipartite graph $G = (U, V, E)$ such that $|U| = 2^n$, $|V| = m$ and
  all $u \in U$ are of degree $\deg(u) \le k$, the following holds.
  There is a constant $C > 0$, depending on the explicitness of $G$,
  such that for all $s \ge C \cdot m \cdot n^C\cdot k^C$ and any
  Boolean function $f \in \calF_n(s/2,t)$ there is a natural
  $m$-independent $k$-determined variable substitutions $\rho$ for the
  formula $\circuit_s(f)$ such that
  \begin{align*}
    \gout_{s, \alpha}(Y) =
    \begin{cases}
      f(\alpha), &\text{if~} C_f(\alpha) \neq \bot,\\
      \oplus_{i \in N(\alpha)} y_i, &\text{otherwise}
    \end{cases}
  \end{align*}
  for all $\alpha \in \set{0,1}^n$ and $\gout_{s, \alpha}$ and $Y$
  as in \cref{def:determined}.  Furthermore, the formula $\restrict{\circuit_s(f)}{\rho}$ is over
  $O\big(t \cdot k + m\big)$ variables.
\end{lemma}

% \pa{Issue with how explicit is defined (see preliminaries) makes this
%   a bit loosey-goosey.  The exact size of the selector circuits
%   depends on $G$ and this governs the hidden constant in the $\Omega$
%   for how large $s$ needs to be.  In particular it is not ``for
%   every'' $s = \Omega(...)$ (which indicates that the hidden constant
%   can be made as small as you want).}
% \kr{Better? At least the statement?}

For the degree lower bound (\cref{thm:main-degree}) we will set $t =
2^n$ and use the trivial $C_f$ which always outputs $\bot$, so the
reader who wishes a simplified version of the lemma can focus on this
special case.

\begin{proof}
  We consider the formula $\circuit_s(f)$ and let the first $m$ gates
  of the formula be denoted by $Y$. We restrict the formula such that
  each gate in $Y$ computes an \emph{or} of two constants. The first
  wire to the gate is fixed to the constant $0$, whereas the second
  wire is only restricted to carry either the constant $0$ or $1$.  In
  the end these will be the only structural variables that are not
  restricted to a constant. In the following we think of the gates $Y$
  as Boolean variables; as $m$ additional input bits to our circuit.

  Further, we restrict another part of the formula such that one part
  of the circuit described by the formula computes the circuit
  $C_f$. Recall that we pretend that the output of $C_f$ is in
  $\set{0,1,\bot}$, but it actually outputs two bits $C_f^1$ and
  $C_f^2$, where $C_f^1(\alpha) = 1$ if and only
  if $C_f^2(\alpha) = f(\alpha)$.

  Finally we also want to hard code the bipartite graph
  $G(\set{0,1}^n, Y, E)$ into our circuit. Since $G$ is very large
  this requires $G$ to be explicit. That is, we require small circuits
  $\sel_1, \ldots, \sel_{m}$, where given any $\alpha \in \set{0,1}^n$, $\sel_i(\alpha)$ is $1$ if and only if the vertex $y_i \in Y$ is a
  neighbor of the vertex $\alpha$.
  %
  By \cref{claim:explicit neigh} these circuits $\sel_i$ are each of size
  \[
    k \cdot \left( \poly(n + \log k) + 2 \log m + 1 \right)
    \le \poly(n, k) \eqperiod
  \]
  %
  The restriction $\rho$ restricts some structural variables such that
  a part of the circuit computes $\sel_1, \ldots, \sel_m$. We connect
  each output of the $\sel_i$ circuit by an \emph{and} gate to the
  negation of $C_f^1$. Denote the resulting circuits by
  $\sel'_1, \ldots, \sel'_m$.  Observe that the circuits $\sel_i'$
  output $0$ whenever $C_f^2(\alpha) = f(\alpha)$ and otherwise output
  $\sel_i$. We think of these circuits as ``selector circuits'' which
  indicate whether on input $\alpha \in \set{0,1}^n$ (to the original
  variables $x_1,\ldots, x_n$ over which the circuit is defined) the
  variable $y_i \in Y$ appears in the constraint for $\alpha$.

  The output of these selector circuits $\sel_i'$ is connected to the
  gate $y_i$ by an \emph{and} gate. All these $m$ \emph{and} gates are
  in turn connected to a circuit computing the \emph{xor} of these
  gates. Finally, to ensure that the circuit computes $f(\alpha)$ on
  inputs $\alpha$ such that $C_f(\alpha) \neq \bot$, we connect
  $C_f^1$ with $C_f^2$ by an \emph{and} gate which is then connceted
  by a \emph{or} gate to the output of the \emph{xor} circuit. This
  completes the description of the restriction on the structure
  variables. A depiction of the resulting circuit can be found in
  \cref{fig:circuit-restricted}.
  \begin{figure}
    \centering
    \includegraphics[width=\linewidth,height=\textheight,keepaspectratio]{fig/circuit-restricted.pdf}
    \caption{A schematic depiction of the formula after hitting it
      with the described restriction.}
    \label{fig:circuit-restricted}
  \end{figure}

  Note that this implements the intended semantics: for each
  input $\alpha \in \set{0,1}^n$ the selector circuits output $1$ on
  some variables $y_i$ which are then \emph{xor}'ed, and the
  restricted circuit outputs
  \begin{align}
    \bigoplus_{i \in N(\alpha)} y_i \eqcomma
  \end{align}
  unless $C_f(\alpha) \neq \bot$, in which case the output of the
  circuit is $f(\alpha)$ and all selector circuits output $0$. We
  require that $s$ is larger than the size of the described circuit
  which is of size $O\big(m \cdot \poly(n, k)\big) + s/2$.

  We have the intended semantics of the circuit and need to ensure the
  furthermore property: that the restricted formula is over few
  variables. First, since the selector circuits $\sel'_i$ are fixed,
  all evaluation variables for these subcircuits can be fixed to
  constants. The same holds for the circuit $C_f$. Similarly, since
  the $y_i$ gate always carries the value of the $y_i$ variable, all
  $2^n \cdot m$ wire variables corresponding to the $Y$ variables can
  be substituted by the corresponding $y_i$ variable and are thus
  restricted away.

  After these restrictions the only evaluation variables left are
  those for the evaluation of the $\oplus$ circuit.  For $\alpha$ such
  that $C_f(\alpha) \neq \bot$, the selector circuits are hard-wired
  to $0$ and in particular the inputs to the $\oplus$ circuit is
  hard-wired to $0$, meaning that these evalation variables can be
  restricted away.

  There remains then only the $O(t \cdot m)$ evaluation variables
  corresponding to the evaluation of the $\oplus$ circuit for inputs
  $\alpha$ such that $C_f(\alpha) = \bot$. Let us, without loss of
  generality, use an \emph{xor}-circuit which iteratively \emph{xor}s
  each variable.  Concretely, let it have subcircuits $\chi_i$ where
  $\chi_1 = \sel'_1 \wedge y_1$ and
  $\chi_{i} = \chi_{i-1} \oplus (\sel'_i \wedge y_i)$ for $i > 1$, and
  $\chi_m$ is the overall output of the $\oplus$ circuit.

  The only observation required is that if the circuit
  $\sel'_i(\alpha) = 0$, then $\chi_i$ gets a $0$ as input from index
  $i$, independent of the value of $y_i$. Hence the output wire
  variable of the circuit $\chi_i$ indexed by the input $\alpha$ can
  be substituted by the output of the circuit $\chi_{i-1}$. Hence for
  each $\alpha$ such that $C_f(\alpha) = \bot$, we can reduce the
  number of free wire variables indexed by $\alpha$ to $O(k)$, as each
  $\oplus$-constraint is over at most $k$ variables.  As $C_f$ outputs
  $\bot$ on at most $t$ inputs, we end up with a restriction leaving
  only a total of $O(t \cdot k + m)$ remaining variables in the
  restricted formula.

  This completes the description of the restriction $\rho$.  The only
  part that remains is to verify that $\rho$ is natural,
  $k$-determined, and $m$-independent.  That $\rho$ is natural is
  immediate -- it does not substitute any structural variable by an
  evaluation variable.  For $k$-determinedness, note that for a fixed
  input $\alpha$ at most $k$ selector circuits output $1$, and thus
  for every gate $u$ the value of $\outwire_\alpha(u)$ as a function
  of $Y$ can be computed by a function over those $k$ variables.
  Finally, each assignment to the remaining structure variables $Y$
  gives a valid circuit and thus $\rho$ is $m$-independent.
\end{proof}

We are ready to prove the degree lower bound, restated here for
convenience.

\MainDegree*

\begin{proof}
  Let $G = (U, V, E)$ be an explicit bipartite graph as in
  \cref{thm:expander}, with $U = \{0,1\}^n$,
  $k = O_{\gamma}\big((n \log r)^{1 + 1/\gamma}\big)$, and
  $|V| \le k^2 r^{1+\gamma}$ for parameters $\gamma > 0$ and
  $r \le 2^n$ to be fixed later. Apply \cref{lem:restriction} with $t=2^n$ along with
  $C_f = \bot$ to obtain, for $s \ge m \cdot \poly(n, k)$, a natural
  $m$-independent $k$-determined variable substitution $\rho$ for
  $\circuit_s(f)$ such that
  $\gout_{s, \alpha}(Y) = \oplus_{i \in N(\alpha)} y_i$. In words, the
  circuit of the restricted formula on input $\alpha$ computes an
  \emph{xor} of the neighborhood of the vertex $\alpha$ of $G$.

  Apply \cref{lem:circuit-to-poly} to $\rho$ to conclude that if there
  is an SoS refutation of $\restrict{\circuit_s(f)}{\rho}$ of degree
  $d$, then there is a degree $d \cdot k$ SoS refutation of the system
  of polynomial equations computing
  \begin{align}
    \nonumber
    \calP_G
    &= \Big\{
    \bigoplus_{i \in N(\alpha)} y_i = f(\alpha)
    :
    \alpha \in \set{0,1}^n
    \Big\}\,
    \cup
    \set{
    y_i^2 = y_i\mid i \in [m]
    }
    \eqperiod
  \end{align}
%
  As the graph $G$ is a strong expander, we can apply
  \cref{thm:sos-xor} to get an SoS degree lower bound of $\Omega(r)$
  for the XOR-CSP instance $\calP_G$ defined over $G$, which in turn
  gives us an $\Omega(r/k)$ degree lower bound for the
  $\restrict{\circuit_s(f)}{\rho}$ formula and hence also for the
  unrestricted formula.

  Let us fix the parameters. We want to choose $r$ as large as
  possible. However, the larger we choose $r$, the larger $m$ may
  become, since \cref{thm:expander} only guarantees that
  $m \le k^2 r^{1+\gamma}$. Let us analyze how large $r$ can be chosen
  in terms of $n$ and $s$.

  Note that $k = \poly_\gamma(n)$, where we use that $r \le 2^n$, and we write
  $\poly_\gamma(n)$ to denote some polynomial in $n$ whose degree and
  coefficients may depend on $\gamma$. Hence we may choose
  \begin{align}
    m = \frac{s}{\poly_{\gamma}(n)} \eqcomma
  \end{align}
  according to the requirement on $s$ in \cref{lem:restriction}.  From
  the guarantees of \cref{thm:expander} we know that
  $r \ge (m/k^2)^{1/(1+\gamma)}$. Substituting $m$ according to the
  previous equation we get that
  \begin{align}
    r
    \ge
    \left(
    \frac{s}{
    k^2 \poly_{\gamma}(n)
    }
    \right)^{\frac{1}{1 + \gamma}}
    =
    \frac{s^{1/(1 + \gamma)}}{\poly_{\gamma}(n)} \eqperiod
  \end{align}
  Hence if we choose $\gamma$ small enough so that
  $\frac{1}{1+\gamma} > 1-\eps/2$ and then require $s$ to be large
  enough such that the final $\poly_{\gamma}(n)$ is at most
  $s^{\eps/2}$, we obtain the claimed lower bound.
\end{proof}

In the following we prove the claimed size lower bound.

\MainSize*

\begin{proof}
  Apply \cref{lem:restriction} with the graphs from
  \cref{thm:expander} as in the proof of \cref{thm:main-degree}. We
  get a natural $m$-independent $k$-determined variable substitution
  $\rho$ and the formula $\restrict{\circuit_s(f)}{\rho}$ over
  $O(t \cdot k + m)$ variables.  To this formula we then apply
  \cref{lem:circuit-to-poly} to obtain a degree lower bound of
  $\Omega(r/k)$, akin to the proof of \cref{thm:main-degree}. By
  setting the parameters as in the aforementioned proof we get the
  same degree lower bound of $\Omega_\eps(s^{1-\eps/3})$ for the
  formula $\restrict{\circuit_s(f)}{\rho}$. As this formula is over
  few variables we can apply \cref{thm:tradeoff} to obtain an SoS size
  lower bound of
  $\exp \Big( \Omega_\eps\big((s^{1- \eps/3} - 3k)^2/(t \cdot k +
  m)\big) \Big)$ for the restricted formula. As variable substitutions
  may only decrease the size of a refutation, the same lower bound
  also holds for the unrestricted formula. We obtain the desired lower
  bound by choosing $s$ large enough such that
  $s^{\eps/3} \ge k = \poly_\eps(n)$ and by recalling that
  $t \ge s \ge m$. 
\end{proof}
